\documentclass[10pt]{article}
\usepackage[utf8]{inputenc}
\usepackage[T1]{fontenc}
\usepackage{amsmath}
\usepackage{amsfonts}
\usepackage{amssymb}
\usepackage[version=4]{mhchem}
\usepackage{stmaryrd}
\usepackage{hyperref}
\hypersetup{colorlinks=true, linkcolor=blue, filecolor=magenta, urlcolor=cyan,}
\urlstyle{same}
\usepackage{graphicx}
\usepackage[export]{adjustbox}
\graphicspath{ {./images/} }
\usepackage{mathrsfs}

\title{FUNCTORS OF ARTIN RINGS( ${ }^{1}$ ) }

\author{BY\\
MICHAEL SCHLESSINGER}
\date{}


%New command to display footnote whose markers will always be hidden
\let\svthefootnote\thefootnote
\newcommand\blfootnotetext[1]{%
  \let\thefootnote\relax\footnote{#1}%
  \addtocounter{footnote}{-1}%
  \let\thefootnote\svthefootnote%
}

%Overriding the \footnotetext command to hide the marker if its value is `0`
\let\svfootnotetext\footnotetext
\renewcommand\footnotetext[2][?]{%
  \if\relax#1\relax%
    \ifnum\value{footnote}=0\blfootnotetext{#2}\else\svfootnotetext{#2}\fi%
  \else%
    \if?#1\ifnum\value{footnote}=0\blfootnotetext{#2}\else\svfootnotetext{#2}\fi%
    \else\svfootnotetext[#1]{#2}\fi%
  \fi
}

\DeclareUnicodeCharacter{21D2}{\ifmmode\Rightarrow\else{$\Rightarrow$}\fi}

\begin{document}
\maketitle


\begin{enumerate}
  \setcounter{enumi}{-1}
  \item Introduction. In the investigation of functors on the category of preschemes, one is led, by Grothendieck [3], to consider the following situation. Let $\Lambda$ be a complete noetherian local ring, $\mu$ its maximal ideal, and $k=\Lambda / \mu$ the residue field. (In most applications $\Lambda$ is $k$ itself, or a ring of Witt vectors.) Let $\boldsymbol{C}$ be the category of Artin local $\Lambda$-algebras with residue field $k$. A covariant functor $F$ from $\boldsymbol{C}$ to Sets is called pro-representable if it has the form
\end{enumerate}

$$
F(A) \cong \operatorname{Hom}_{\text {local } \Lambda \text {-alg. }}(R, A), \quad A \in C,
$$

where $R$ is a complete local $\Lambda$-algebra such that $R / \boldsymbol{m}^{n}$ is in $\boldsymbol{C}$, all $n$. ( $\boldsymbol{m}$ is the maximal ideal in $R$.)

In many cases of interest, $F$ is not pro-representable, but at least one may find an $R$ and a morphism $\operatorname{Hom}(R, \cdot) \rightarrow F$ of functors such that $\operatorname{Hom}(R, A) \rightarrow F(A)$ is surjective for all $A$ in $C$. If $R$ is chosen suitably "minimal" then $R$ is called a "hull" of $F ; R$ is then unique up to noncanonical isomorphism. Theorem 2.11, §2, gives a criterion for $F$ to have a hull, and also a simple criterion for pro-representability which avoids the use of Grothendieck's techniques of nonflat descent [3], in some cases. Grothendieck's program is carried out by Levelt in [4]. §3 contains a few geometric applications of these results.\\[0pt]
To avoid awkward terminology, I have used the word "pro-representable" in a more restrictive sense than Grothendieck [3] has. He considers the category of $\Lambda$-algebras of finite length and allows $R$ to be a projective limit of such rings.\\
The methods of this paper are a simple extension of those used by David Mumford in a proof (unpublished) of the existence of formal moduli for polarized Abelian varieties. I am indebted to Mumford and to John Tate for many valuable suggestions.

\begin{enumerate}
  \item The category $\boldsymbol{C}_{\Lambda}$. Let $\Lambda$ be a complete noetherian local ring, with maximal ideal $\boldsymbol{\mu}$ and residue field $k=\Lambda / \boldsymbol{\mu}$. We define $\boldsymbol{C}=\boldsymbol{C}_{\Lambda}$ to be the category of Artinian local $\Lambda$-algebras having residue field $k$. (That is, the "structure morphism" $\Lambda \rightarrow A$ of such a ring $A$ induces a trivial extension of residue fields.) Morphisms in $C$ are local homomorphisms of $\Lambda$-algebras.
\end{enumerate}

\footnotetext{\href{http://Received.by}{Received.by} the editors March 8, 1966.\\
( ${ }^{1}$ ) The contents of this paper form part of the author's 1964 Harvard Ph.D. Thesis, which was directed by John Tate. This research was supported in part by a grant from the Air Force Office of Scientific Research.
}Let $\hat{\boldsymbol{C}}=\hat{\boldsymbol{C}}_{\Delta}$ be the category of complete noetherian local $\Lambda$-algebras $A$ for which $A / \boldsymbol{m}^{n}$ is in $\boldsymbol{C}$, all $n$. Notice that $\boldsymbol{C}$ is a full subcategory of $\hat{\boldsymbol{C}}$.

If $p: A \rightarrow B, q: C \rightarrow B$ are morphisms in $C$, let $A \times_{B} C$ denote the ring (in $C$ ) consisting of all pairs $(a, c)$ with $a \in A, c \in C$, for which $p a=q c$, with coordinatewise multiplication and addition.

For any $A$ in $\hat{\boldsymbol{C}}$, we denote by $t_{A / \Lambda}^{*}$, or just $t_{A}^{*}$, the "Zariski cotangent space" of $A$ over $\Lambda$ :


\begin{equation*}
t_{A}^{*}=m /\left(m^{2}+\mu A\right) \tag{1.0}
\end{equation*}


where $\boldsymbol{m}$ is the maximal ideal of $A$. A simple calculation shows that the dual vector space, denoted by $t_{A}$, may be identified with $\operatorname{Der}_{\Lambda}(A, k)$, the space of $\Lambda$ linear derivations of $A$ into $k$.

Lemma 1.1. A morphism $B \rightarrow A$ in $\hat{\boldsymbol{C}}$ is surjective if and only if the induced map from $t_{B}^{*}$ to $t_{A}^{*}$ is surjective.

Proof. First of all, any $A$ in $\hat{C}$ is generated, as $\Lambda$ module, by the image of $\Lambda$ in $A$ and the maximal ideal $\boldsymbol{m}$ of $A$. (For $A$ and $\Lambda$ have the same residue field $k$.) Thus the induced map from $\mu / \mu^{2}$ to $\mu A /\left(\boldsymbol{m}^{2} \cap \mu A\right)$ is a surjection. If $B \rightarrow A$ is a morphism in $\hat{\boldsymbol{C}}$, then denoting the maximal ideal of $B$ by $\boldsymbol{n}$, we get a commutative diagram with exact rows:\\
\includegraphics[max width=\textwidth, alt={}, center]{133f4430-1007-4a22-9674-9093b9475bf3-02_188_762_1277_452}\\
in which the left-hand arrow is a surjection. If the right-hand arrow is also a surjection, then the middle arrow is a surjection, so that the induced map on the graded rings is a surjection. From this it follows that $B \rightarrow A$ is a surjection [1, §2, No. 8, Theorem 1].

Conversely, if $B \rightarrow A$ is a surjection, then the induced map on cotangent spaces is obviously surjective.

Let $p: B \rightarrow A$ be a surjection in $C$.\\
Definition 1.2. $p$ is a small extension if kernel $p$ is a nonzero principal ideal (t) such that $\boldsymbol{m} \boldsymbol{t}=(\mathbf{0})$, where $\boldsymbol{m}$ is the maximal ideal of $\boldsymbol{B}$.

Definition 1.3. $p$ is essential if for any morphism $q: C \rightarrow B$ in $C$ such that $p q$ is surjective, it follows that $q$ is surjective.

From Lemma 1.1 we obtain easily\\
Lemma 1.4. Let $p: B \rightarrow A$ be a surjection in $C$. Then

\begin{itemize}
  \item[(i)] $p$ is essential if and only if the induced map $p_{*}: t_{B}^{*} \rightarrow t_{A}^{*}$ is an isomorphism.
  \item[(ii)] If $p$ is a small extension, then $p$ is not essential if and only if $p$ has a section $s: A \rightarrow B$, with $p s=1_{A}$.
\end{itemize}

Proof. (i) If $p_{*}$ is an isomorphism, then by Lemma 1.1, $p$ is essential. Conversely let $\tilde{t}_{1}, \ldots, \tilde{t}_{r}$ be a basis of $t_{A}^{*}$, and lift the $\tilde{t}_{i}$ back to elements $t_{i}$ in $B$. Set

$$
C=\Lambda\left[t_{1}, \ldots, t_{r}\right] \subseteq B .
$$

Then $p$ induces a surjection from $C$ to $A$, so if $p$ is essential, $C=B$. But then $\operatorname{dim}_{k} t_{B}^{*} \leqq r=\operatorname{dim}_{k} t_{A}^{*}$, so $t_{B}^{*} \cong t_{A}^{*}$.\\
(ii) If $p$ has a section $s$, then $s$ is not surjective, so $p$ is not essential. If $p$ is not essential, then the subring $C$ constructed above is a proper subring of $B$, and hence is isomorphic to $A$, since length $(B)=$ length $(A)+1$. The isomorphism $C \cong A$ yields the section.\\
2. Functors on $\boldsymbol{C}$. We shall consider only ccvariant functors $F$, from $\boldsymbol{C}$ to Sets, such that $F(k)$ contains just one element. By a couple for $F$ we mean a pair $(A, \xi)$ where $A \in C$ and $\xi \in F(A)$. A morphism of couples $u:(A, \xi) \rightarrow\left(A^{\prime}, \xi^{\prime}\right)$ is a morphism $u: A \rightarrow A^{\prime}$ in $\boldsymbol{C}$ such that $F(u)(\xi)=\xi^{\prime}$. If we extend $F$ to $\hat{\boldsymbol{C}}$ by the formula $\hat{F}(A)=$ proj $\operatorname{Lim} F\left(A / \boldsymbol{m}^{n}\right)$ we may speak analogously of pro-couples and morphisms of pro-couples.

For any ring $R$ in $\hat{\boldsymbol{C}}$, we set $h_{R}(A)=\operatorname{Hom}(R, A)$ to define a functor $h_{R}$ on $C$. Now if $F$ is any functor on $C$, and $R$ is in $\hat{\boldsymbol{C}}$, we have a canonical isomorphism

$$
\hat{F}(R) \xrightarrow{\sim} \operatorname{Hom}\left(h_{R}, F\right) .
$$

Namely, let $\xi=$ proj $\operatorname{Lim} \xi_{n}$ be in $\hat{F}(R)$. Then each $u: R \rightarrow A$ factors through $u_{n}: R / \boldsymbol{m}^{n} \rightarrow \mathcal{A}$ for some $n$, and we assign to $u \in h_{R}(A)$ the element $F\left(u_{n}\right)\left(\xi_{n}\right)$ of $F(A)$. This sets up the isomorphism. We therefore say that a pro-couple $(R, \xi)$ for $F$ pro-represents $F$ if the morphism $h_{R} \rightarrow F$ induced by $\xi$ is an isomorphism.\\
(2.1) Relation to global functors. Let $G$ be a contravariant functor on the category of preschemes over Spec $\Lambda$, and pick a fixed $e \in G(\operatorname{Spec} k)$. For $A$ in $C$, let $F(A) \subseteq G(\operatorname{Spec} A)$ be the set of those $\xi \in G(\operatorname{Spec} A)$ such that $G(i)(\xi)=e$ where $i$ is the inclusion of $\operatorname{Spec} k$ in $\operatorname{Spec} A$. If $G$ is represented by a prescheme $X$, then $e$ determines a $k$-rational point $x \in X$, and it is then clear that $F(A)$ is isomorphic to $\operatorname{Hom}_{\Lambda}\left(\mathfrak{O}_{X, x}, A\right)$. Thus the completion of $\mathfrak{O}_{X, x}$ pro-represents $F$.

Unfortunately, many interesting functors, for example some "formal moduli" functors (§3.7), are not pro-representable. However, one can still look for a "universal object" in some sense, for example in the sense of Definition 2.7 below.

Definition 2.2. A morphism $F \rightarrow G$ of functors is smooth if for any surjection $B \rightarrow A$ in $C$, the morphism


\begin{equation*}
F(B) \rightarrow F(A) \times_{G(A)} G(B) \tag{*}
\end{equation*}


is surjective.\\
Part (i) of the sorités below will perhaps motivate this definition.\\
Remarks. (2.3) It is enough to check surjectivity in (*) for small extensions $\boldsymbol{B} \rightarrow \boldsymbol{A}$.\\
(2.4) If $F \rightarrow G$ is smooth, then $\hat{F} \rightarrow \hat{G}$ is surjective, in the sense that $\hat{F}(A) \rightarrow \hat{G}(A)$ is surjective for all $A$ in $\hat{\boldsymbol{C}}$ (consider the successive quotients $A / \boldsymbol{m}^{n}, n=1,2, \ldots$ ).

Proposition 2.5. (i) Let $R \rightarrow S$ be a morphism in $\hat{\boldsymbol{C}}$. Then $h_{S} \rightarrow h_{R}$ is smooth if and only if $S$ is a power series ring over $R$.\\
(ii) If $F \rightarrow G$ and $G \rightarrow H$ are smooth morphisms of functors, then the composition $F \rightarrow H$ is smooth.\\
(iii) If $u: F \rightarrow G$ and $v: G \rightarrow H$ are morphisms of functors such that $u$ is surjective and $v u$ is smooth, then $v$ is smooth.\\
(iv) If $F \rightarrow G$ and $H \rightarrow G$ are morphisms of functors such that $F \rightarrow G$ is smooth, then $F \times_{G} H \rightarrow H$ is smooth.

Proof. (i) This is more or less well known (see [3, Theorem 3.1]), but we give a proof for the sake of completeness. Suppose $h_{S} \rightarrow h_{R}$ is smooth. Let $r$ (resp. $s$ ) be the maximal ideal in $R$ (resp. $S$ ), and pick $x_{1}, \ldots, x_{n}$ in $S$ which induce a basis of $t_{S / R}^{*}=s /\left(s^{2}+r S\right)$. If we set $T=R\left[\left[X_{1}, \ldots, X_{n}\right]\right]$ and denote the maximal ideal of $T$ by $t$, we get a morphism $u_{1}: S \rightarrow T /\left(t^{2}+r T\right)$ of local $R$ algebras, obtained by mapping $x_{i}$ on the residue class of $X_{i}$. By smoothness $u_{1}$ lifts to $u_{2}: S \rightarrow T / t^{2}$, thence to $u_{3}: S \rightarrow T / t^{3}, \ldots$ etc. Thus we get a $u: S \rightarrow T$ which induces an isomorphism of $t_{S / R}^{*}$ with $t_{T / R}^{*}$ (by choice of $u_{1}$ ) so that $u$ is a surjection (1.1). Furthermore, if we choose $y_{i} \in S$ such that $u y_{i}=X_{i}$, we can set $v X_{i}=y_{i}$ and produce a local morphism $v: T \rightarrow S$ of $R$ algebras such that $u v=1_{T}$; in particular $v$ is an injection. Clearly $v$ induces a bijection on the cotangent spaces, so $v$ is also a surjection (1.1). Hence $v$ is an isomorphism of $T=R\left[\left[X_{1}, \ldots, X_{n}\right]\right]$ with $S$.

Conversely, if $S$ is a power series ring over $R$, then it is obvious that $h_{S} \rightarrow h_{R}$ is smooth.

The proofs of (ii), (iii), (iv) are completely formal and are left to the reader.\\
(2.6) Notation. Let $k[\varepsilon]$, where $\varepsilon^{2}=0$, denote the ring of dual numbers over $k$. For any functor $F$, the set $F(k[\varepsilon])$ is called the tangent space to $F$, and is denoted by $t_{F}$. It is easy to see that if $F=h_{R}$, then there is a canonical isomorphism $t_{F} \cong t_{R}$ :

$$
t_{R} \cong \operatorname{Hom}_{\Lambda}(R, k[\varepsilon]) .
$$

Usually $t_{F}$ will have an intrinsic vector space structure (Lemma 2.10).\\
Definition 2.7. A pro-couple $(R, \xi)$ for a functor $F$ is called a pro-representable hull of $F$, or just a hull of $F$, if the induced map $h_{R} \rightarrow F$ is smooth (2.2), and if in addition the induced map $t_{R} \rightarrow t_{F}$ of tangent spaces is a bijection.\\
(2.8) Notice that if $(R, \xi)$ pro-represents $F$ then $(R, \xi)$ is a hull of $F$. In this case $(R, \xi)$ is unique up to canonical isomorphism. In general we have only noncanonical isomorphism:

Proposition 2.9. Let $(R, \xi)$ and $\left(R^{\prime}, \xi^{\prime}\right)$ be hulls of $F$. Then there exists an isomorphism $u: R \rightarrow R^{\prime}$ such that $F(u)(\xi)=\xi^{\prime}$.

Proof. By (2.4) we have morphisms $u:(R, \xi) \rightarrow\left(R^{\prime}, \xi^{\prime}\right)$ and $u^{\prime}:\left(R^{\prime}, \xi^{\prime}\right) \rightarrow(R, \xi)$, both inducing an isomorphism on tangent spaces, by the definition of hull. Thus\\
$u^{\prime} u$ say, induces an isomorphism on $t_{R}^{*}$, so that $u^{\prime} u$ is a surjective endomorphism of $R$, by Lemma 1.1. But an easy argument, which we leave to the reader, shows that a surjective endomorphism of any noetherian ring is an isomorphism. Thus $u^{\prime} u$ and $u u^{\prime}$ are isomorphisms and we are done.

Remark 2.10. Let $(R, \xi)$ be a hull of $F$. Then $R$ is a power series ring over $\Lambda$ if and only if $F$ transforms surjections $B \rightarrow A$ in $C$ into surjections $F(B) \rightarrow F(A)$. In fact the stated condition on $F$ is equivalent to the smoothness of the natural morphism $F \rightarrow h_{\Lambda}$. By applying (2.6), (ii) and (iii) to the diagram\\
\includegraphics[max width=\textwidth, alt={}, center]{133f4430-1007-4a22-9674-9093b9475bf3-05_171_268_708_697}\\
we conclude that $h_{R} \rightarrow h_{\Lambda}$ is smooth if and only if $F \rightarrow h_{\Lambda}$ is. Now use 2.5 (i).\\
Lemma 2.10. Suppose $F$ is a functor such that

$$
F\left(k[V] \times_{k} k[W]\right) \xrightarrow{\sim} F(k[V]) \times F(k[W])
$$

for vector spaces $V$ and $W$ over $k$, where $k[V]$ denotes the ring $k \oplus V$ of $\boldsymbol{C}$ in which $V$ is a square zero ideal. Then $F(k[V])$, and in particular $t_{F}=F(k[\varepsilon])$, has a canonical vector space structure, such that $F(k[V]) \cong t_{F} \otimes V$.

Proof. $k[V]$ is in fact a "vector space object" in the category $\hat{\boldsymbol{C}}$ (in which $k$ is the final object), for we have a canonical isomorphism

$$
\operatorname{Hom}(A, k[V]) \cong \operatorname{Der}_{\Lambda}(A, V), \quad A \in \hat{C} .
$$

The addition map $k[V] \times_{k} k[V] \rightarrow k[V]$ is given by $(x, 0) \mapsto x,(0, x) \mapsto x(x \in V)$, and scalar multiplication by $a \in k$ is given by the endomorphism $x \mapsto a x(x \in V)$ of $k[V]$. Thus if $F$ commutes with the necessary products, $F(k[V])$ gets a vector space structure. Finally, we identify $V$ with $\operatorname{Hom}(k[\varepsilon], k[V])$ to get a map

$$
t_{F} \otimes V \rightarrow F(k[V])
$$

which is an isomorphism since $k[V]$ is isomorphic to the product of $r=\operatorname{dim}_{k} V$ copies of $k[\varepsilon]$.

Theorem 2.11. Let $F$ be a functor from $\boldsymbol{C}$ to Sets such that $F(k)=(e)$ (=one point). Let $A^{\prime} \rightarrow A$ and $A^{\prime \prime} \rightarrow A$ be morphisms in $C$, and consider the map


\begin{equation*}
F\left(A^{\prime} \times_{A} A^{\prime \prime}\right) \rightarrow F\left(A^{\prime}\right) \times_{F(A)} F\left(A^{\prime \prime}\right) . \tag{2.12}
\end{equation*}


Then

\begin{itemize}
  \item[(1)] $F$ has a hull if and only if $F$ has properties $\left(\mathrm{H}_{1}\right),\left(\mathrm{H}_{2}\right),\left(\mathrm{H}_{3}\right)$ below:
  \item[] $\left(\mathrm{H}_{1}\right)(2.12)$ is a surjection whenever $A^{\prime \prime} \rightarrow A$ is a small extension (1.2).
  \item[] $\left(\mathrm{H}_{2}\right)(2.12)$ is a bijection when $A=k, A^{\prime \prime}=k[\varepsilon]$.
  \item[] $\left(\mathrm{H}_{3}\right) \operatorname{dim}_{k}\left(t_{F}\right)<\infty$.
\end{itemize}

(2) $F$ is pro-representable if and only if $F$ has the additional property $\left(\mathrm{H}_{4}\right)$ :


\begin{equation*}
F\left(A^{\prime} \times{ }_{A} A^{\prime}\right) \xrightarrow{\sim} F\left(A^{\prime}\right) \times_{F(A)} F\left(A^{\prime}\right) \tag{4}
\end{equation*}


for any small extension $A^{\prime} \rightarrow A$.\\
Notice that if $F$ is isomorphic to some $h_{R}$, then (2.12) is an isomorphism for any morphisms $A^{\prime} \rightarrow A, A^{\prime \prime} \rightarrow A$; that is, the four conditions are trivially necessary for pro-representability.

Remarks. (2.13) ( $\mathrm{H}_{2}$ ) implies that $t_{F}$ is a vector space by Lemma 2.10. In fact, by induction on $\operatorname{dim}_{k} W$ we conclude from $\left(\mathrm{H}_{2}\right)$ that (2.12) is an isomorphism when $A=k, A^{\prime \prime}=k[W]$; in particular the hypotheses of 2.10 are satisfied.\\
(2.14) By induction on length $A^{\prime \prime}$-length $A$ it follows from ( $\mathrm{H}_{1}$ ) that (2.12) is a surjection for any surjection $A^{\prime \prime} \rightarrow A$.\\
(2.15) Condition $\left(\mathrm{H}_{4}\right)$ may be usefully viewed as follows. For each $A$ in $C$, and each ideal $I$ in $A$ such that $m_{A} \cdot I=(0)$, we have an isomorphism


\begin{equation*}
A \times_{A / I} A \xrightarrow{\sim} A \times_{k} k[I], \tag{2.16}
\end{equation*}


induced by the map $(x, y) \mapsto\left(x, x_{0}+y-x\right)$, where $x$ and $y$ are in $A$ and $x_{0}$ is the $k$ residue of $x$. Now, given a small extension $p: A^{\prime} \rightarrow A$ with kernel $I$, we get by $\left(\mathrm{H}_{2}\right)$ and (2.16) a map


\begin{equation*}
F\left(A^{\prime}\right) \times\left(t_{F} \otimes I\right) \rightarrow F\left(A^{\prime}\right) \times_{F(A)} F\left(A^{\prime}\right) \tag{2.17}
\end{equation*}


which is easily seen to determine, for each $\eta \in F(A)$, a group action of $t_{F} \otimes I$ on the subset $F(p)^{-1}(\eta)$ of $F\left(A^{\prime}\right)$ (provided that subset is not empty). ( $\mathrm{H}_{1}$ ) implies that this action is "transitive," while $\left(\mathrm{H}_{4}\right)$ is precisely the condition that this action makes $F(p)^{-1}(\eta)$ a (formally) principal homogeneous space under $t_{F} \otimes I$. Thus, in the presence of conditions $\left(\mathrm{H}_{1}\right),\left(\mathrm{H}_{2}\right),\left(\mathrm{H}_{3}\right)$, it is the existence of "fixed points" of $t_{F} \otimes I$ acting on $F(p)^{-1}(\eta)$ which obstruct the pro-representability of $F$. In many applications, where the elements of $F(A)$ are isomorphism classes of geometric objects, the existence of such a fixed point $\eta^{\prime} \in F(p)^{-1}(\eta)$ is equivalent to the existence of an automorphism of an object $y$ in the class of $\eta$ which cannot be extended to an automorphism of any (or some) object $y^{\prime}$ in the class of $\eta^{\prime}$.

Proof of 2.11. (1) Suppose $F$ satisfies conditions $\left(\mathrm{H}_{1}\right),\left(\mathrm{H}_{2}\right),\left(\mathrm{H}_{3}\right)$. Let $t_{1}, \ldots, t_{r}$ be a dual basis of $t_{F}$, put $S=\Lambda\left[\left[T_{1}, \ldots, T_{r}\right]\right]$, and let $\boldsymbol{n}$ be the maximal ideal of $S$. $R$ will be constructed as the projective limit of successive quotients of $S$. To begin, let $R_{2}=S /\left(n^{2}+\mu S\right) \cong k[\varepsilon] \times_{k} \cdots \times_{k} k[\varepsilon]$ ( $r$ times). By $\left(\mathrm{H}_{2}\right)$ there exists $\xi_{2} \in F\left(R_{2}\right)$ which induces a bijection between $t_{R_{2}}\left(\cong \operatorname{Hom}\left(R_{2}, k[\varepsilon]\right)\right)$ and $t_{F}$. Suppose we have found $\left(R_{q}, \xi_{q}\right)$, where $R_{q}=S / J_{q}$. We seek an ideal $J_{q+1}$ in $S$, minimal among those ideals $J$ in $S$ satisfying the conditions (a) $n J_{q} \subseteq J \subseteq J_{q}$, (b) $\xi_{q}$ lifts to $S / J$. Since the set $\mathscr{S}$ of such ideals corresponds to a certain collection of vector subspaces of $J_{q} /\left(n J_{q}\right)$, it suffices to show that $\mathscr{S}$ is stable under pairwise intersection. But if\\
$J$ and $K$ are in $\mathscr{S}$, we may enlarge $J$, say, so that $J+K=J_{q}$, without changing the intersection $J \cap K$. Then

$$
S / J \times_{S / J_{q}} S / K \cong S /(J \cap K)
$$

so that by $\left(\mathrm{H}_{1}\right)$ (see (2.14)) we may conclude that $J \cap K$ is in $\mathscr{S}$. Let $J_{q+1}$ be the intersection of the members of $\mathscr{S}$, put $R_{q+1}=S / J_{q+1}$, and pick any $\xi_{q+1} \in F\left(R_{q+1}\right)$ which projects onto $\xi_{q} \in F\left(R_{q}\right)$.

Now let $J$ be the intersection of all the $J_{q}$ 's $(q=2,3, \ldots)$ and let $R=S / J$. Since $n^{q} \subseteq J_{q}$, the $J_{q} / J$ form a base for the topology in $R$, so that $R=$ proj $\operatorname{Lim} S / J_{q}$, and it is legitimate to set $\xi=$ proj $\operatorname{Lim} \xi_{q} \in \hat{F}(R)$. Notice that $t_{F} \cong t_{R}$, by choice of $R_{2}$.

We claim now that $h_{R} \rightarrow F$ is smooth. Let $p:\left(A^{\prime}, \eta^{\prime}\right) \rightarrow(A, \eta)$ be a morphism of couples, where $p$ is a small extension, $A=A^{\prime} / I$, and let $u:(R, \xi) \rightarrow(A, \eta)$ be a given morphism. We have to lift $u$ to a morphism $(R, \xi) \rightarrow\left(A^{\prime}, \eta^{\prime}\right)$. For this it suffices to find a $u^{\prime}: R \rightarrow A^{\prime}$ such that $p u^{\prime}=u$. In fact, we have a transitive action of $t_{F} \otimes I$ on $F(p)^{-1}(\eta)$ (resp. $h_{R}(p)^{-1}(\eta)$ ) by (2.15); thus, given such a $u^{\prime}$, there exists $\sigma \in t_{F} \otimes I$ such that $\left[F\left(u^{\prime}\right)(\xi)\right]^{\sigma}=\eta^{\prime}$, so that $v^{\prime}=\left(u^{\prime}\right)^{\sigma}$ will satisfy $F\left(v^{\prime}\right)(\xi)=\eta^{\prime}$, $p v^{\prime}=u$.

Now $u$ factors as $(R, \xi) \rightarrow\left(R_{q}, \xi_{q}\right) \rightarrow(A, \eta)$ for some $q$. Thus it suffices to complete the diagram\\
\includegraphics[max width=\textwidth, alt={}, center]{133f4430-1007-4a22-9674-9093b9475bf3-07_216_313_1247_671}\\
or equivalently, the diagram\\
\includegraphics[max width=\textwidth, alt={}, center]{133f4430-1007-4a22-9674-9093b9475bf3-07_223_596_1595_519}\\
where $w$ has been chosen so as to make the square commute. If the small extension $p r_{1}$ has a section, then $v$ obviously exists. Otherwise, by 1.4 (ii), $p r_{1}$ is essential, so $w$ is a surjection. By $\left(\mathrm{H}_{1}\right)$, applied to the projections of $R_{q} \times_{A} A^{\prime}$ on its factors, $\xi_{q} \in F\left(R_{q}\right)$ lifts back to $R_{q} \times_{A} A^{\prime}$, so ker $w \supseteq J_{q+1}$, by choice of $J_{q+1}$. Thus $w$ factors through $S / J_{q+1}=R_{q+1}$, and $v$ exists. This completes the proof that $(R, \xi)$ is a hull of $F$.

Conversely, suppose that a pro-couple $(R, \xi)$ is a hull of $F$. To verify $\left(\mathrm{H}_{1}\right)$, let $p^{\prime}:\left(A^{\prime}, \eta^{\prime}\right) \rightarrow(A, \eta)$ and $p^{\prime \prime}:\left(A^{\prime \prime}, \eta^{\prime \prime}\right) \rightarrow(A, \eta)$ be morphisms of couples, where $p^{\prime \prime}$\\
is a surjection. Since $h_{R} \rightarrow F$ is surjective, there exists a $u^{\prime}:(R, \xi) \rightarrow\left(A^{\prime}, \eta^{\prime}\right)$, and hence by smoothness applied to $p^{\prime \prime}$, there exists $u^{\prime \prime}:(R, \xi) \rightarrow\left(A^{\prime \prime}, \eta^{\prime \prime}\right)$ rendering the following diagram commutative:\\
\includegraphics[max width=\textwidth, alt={}, center]{133f4430-1007-4a22-9674-9093b9475bf3-08_421_627_507_519}\\
Therefore $\zeta=F\left(u^{\prime} \times u^{\prime \prime}\right)(\xi)$ projects onto $\eta^{\prime}$ and $\eta^{\prime \prime}$, so that $\left(\mathrm{H}_{1}\right)$ is satisfied.\\
Now suppose $(A, \eta)=(k, e)$, and $A^{\prime \prime}=k[\varepsilon]$. If $\zeta_{1}$ and $\zeta_{2}$ in $F\left(A^{\prime} \times_{k} k[\varepsilon]\right)$ have the same projections $\eta^{\prime}$ and $\eta^{\prime \prime}$ on the factors, then choosing $u^{\prime}$ as above we get morphisms

$$
u^{\prime} \times u_{i}:(R, \xi) \rightarrow\left(A^{\prime} \times_{k} k[\varepsilon], \zeta_{i}\right), \quad i=1,2,
$$

by smoothness applied to the projection of $A^{\prime} \times_{k} k[\varepsilon]$ on $A^{\prime}$. Since $t_{F} \cong t_{R}$ we have $u_{1}=u_{2}$, so that $\zeta_{1}=\zeta_{2}$, which proves $\left(\mathrm{H}_{2}\right)$. The isomorphism $t_{R} \cong t_{F}$ also proves $\left(\mathrm{H}_{3}\right)$.\\
(2) Suppose now that $F$ satisfies conditions $\left(\mathrm{H}_{1}\right)$ through $\left(\mathrm{H}_{4}\right)$. By part (1) we know that $F$ has a hull $(R, \xi)$. We shall prove that $h_{R}(A) \xrightarrow{\sim} F(A)$ by induction on length $A$. Consider a small extension $p: A^{\prime} \rightarrow A=A^{\prime} \mid I$, and assume that $h_{R}(A) \xrightarrow{\sim} F(A)$. For each $\eta \in F(A), h_{R}(p)^{-1}(\eta)$ and $F(p)^{-1}(\eta)$ are both formally principal homogeneous spaces under $t_{F} \otimes I(2.15)$; since $h_{R}\left(A^{\prime}\right)$ maps onto $F\left(A^{\prime}\right)$, we have $h_{R}\left(A^{\prime}\right) \xrightarrow{\sim} F\left(A^{\prime}\right)$, which proves the induction step.

The necessity of the four conditions has already been noted.\\
3. Examples.\\
(3.1) The Picard functor. If $X$ is a prescheme, we define $\operatorname{Pic}(X)=H^{1}\left(X, \mathfrak{D}_{X}^{*}\right)$, the group of isomorphism classes of invertible (i.e., locally free of rank one) sheaves on $X$. Recall that the group of automorphisms of an invertible sheaf is canonically isomorphic to $H^{0}\left(X, \mathfrak{O}_{X}^{*}\right)$.

Now suppose $X$ is a prescheme over Spec $\Lambda$. We let $X_{A}$ abbreviate $X \times_{\operatorname{Spec} \Lambda} \operatorname{Spec} A$ for $A$ in $C$, and set $X_{0}=X_{k}$. If $\eta$ (resp. $L$ ) is an element of $\operatorname{Pic}\left(X_{A}\right)$ (resp. an invertible sheaf on $X_{A}$ ) and $A \rightarrow B$ is a morphism in $C$, let $\eta \otimes_{A} B$ (resp. $L \otimes_{A} B$ ) denote the induced element of $\operatorname{Pic}\left(X_{B}\right)$ (resp. induced invertible sheaf on $X_{B}$ ). Let $\xi_{0}$ be an element of $\operatorname{Pic}\left(X_{0}\right)$ fixed once and for all in this discussion, and let\\
$\boldsymbol{P}(A)$ be the subset of $\operatorname{Pic}\left(X_{A}\right)$ consisting of those $\eta$ such that $\eta \otimes_{A} k=\xi_{0}$. We claim that $P$ is pro-representable under suitable conditions, namely:

Proposition 3.2. Assume

\begin{itemize}
  \item[(i)] $X$ is flat over $\Lambda$,
  \item[(ii)] $A \xrightarrow{\sim} H^{\circ}\left(X_{A}, \mathfrak{O}_{X_{A}}\right)$ for each $A \in C$,
  \item[(iii)] $\operatorname{dim}_{k} H^{1}\left(X_{0}, \mathfrak{O}_{X_{0}}\right)<\infty$.
\end{itemize}

Then $\boldsymbol{P}$ is pro-representable by a pro-couple $(R, \xi)$; furthermore $t_{R} \cong H^{1}\left(X_{0}, \mathfrak{O}_{X_{0}}\right)$.\\
Notice that condition (ii) is equivalent to the condition $k \longrightarrow H^{0}\left(X_{0}, \mathfrak{D}_{X_{0}}\right)$, in view of (i). In fact, by flatness, the functor $M \mapsto T(M)=H^{0}\left(X, \mathfrak{O}_{X} \otimes M\right)$ of $\Lambda$ modules is left exact. A standard five lemma type of argument then shows that the natural map $M \rightarrow T(M)$ is an isomorphism for all $M$ of finite length.

For the proof of 3.2 we need two simple lemmas on flatness.\\
Lemma 3.3. Let $A$ be a ring, $J$ a nilpotent ideal in $A$, and $u: M \rightarrow N$ a homomorphism of $A$ modules, with $N$ flat over $A$. If $\bar{u}: M / J M \rightarrow N / J N$ is an isomorphism, then $u$ is an isomorphism.

Proof. Let $K=$ coker $u$ and tensor the exact sequence «

$$
M \rightarrow N \rightarrow K \rightarrow 0
$$

with $A / J$. Then we find $K / J K=0$, which implies $K=0$, since $J$ is nilpotent. Thus, if $K^{\prime}=$ ker $u$, we get an exact sequence

$$
0 \rightarrow K^{\prime} / J K^{\prime} \rightarrow M / J M \rightarrow N / J N \rightarrow 0
$$

by the flatness of $N$. Hence $K^{\prime}=0$, so that $u$ is an isomorphism.\\
Lemma 3.4. Consider a commutative diagram\\
\includegraphics[max width=\textwidth, alt={}, center]{133f4430-1007-4a22-9674-9093b9475bf3-09_419_389_1601_651}\\
of compatible ring and module homomorphisms, where $B=A^{\prime} \times_{A} A^{\prime \prime}, N=M^{\prime} \times_{M} M^{\prime \prime}$ and $M^{\prime}$ (resp. $M^{\prime \prime}$ ) is a flat $A^{\prime}$ (resp. $A^{\prime \prime}$ ) module. Suppose

\begin{itemize}
  \item[(i)] $A^{\prime \prime} / J \xrightarrow{\sim} A$, where $J$ is a nilpotent ideal in $A^{\prime \prime}$,
  \item[(ii)] $u^{\prime}$ (resp. $u^{\prime \prime}$ ) induces $M^{\prime} \otimes_{A^{\prime}} A \xrightarrow{\sim} M\left(\right.$ resp. $\left.M^{\prime \prime} \otimes_{A^{\prime \prime}} A \xrightarrow{\sim} M\right)$.
\end{itemize}

Then $N$ is flat over $B$ and $p^{\prime}$ (resp. $p^{\prime \prime}$ ) induces $N \otimes_{B} A^{\prime} \xrightarrow{\sim} M^{\prime} \quad$ (resp. $\left.N \otimes_{B} A^{\prime \prime} \xrightarrow{\sim} M^{\prime \prime}\right)$.

Proof. We shall consider only the case where $M^{\prime}$ is actually a free $A^{\prime}$ module. (This case actually suffices for our purposes, since a simple application of Lemma 3.3 shows that a flat module over an Artin local ring is free.) Choose a basis $\left(x_{i}^{\prime}\right)_{i \in I}$ for $M^{\prime}$. Then by (ii) we find that $M$ is the free module on generators $u^{\prime}\left(x_{i}^{\prime}\right)$. Choosing $x_{i}^{\prime \prime} \in M^{\prime \prime}$ such that $u^{\prime \prime}\left(x_{i}^{\prime \prime}\right)=u^{\prime}\left(x_{i}^{\prime}\right)$, we get a map $\sum A^{\prime \prime} x_{i}^{\prime \prime} \rightarrow M^{\prime \prime}$ of $A^{\prime \prime}$ modules, whose reduction modulo the ideal $J$ is an isomorphism. Therefore $M^{\prime \prime}$ is free on generators $x_{i}^{\prime \prime}$ (Lemma 3.3) and it follows easily that $N=M^{\prime} \times_{M} M^{\prime \prime}$ is free on generators $x_{i}^{\prime} \times x_{i}^{\prime \prime}$, and that the projections on the factors induce isomorphisms

$$
N \otimes_{B} A^{\prime} \xrightarrow{\sim} M^{\prime}, \quad N \otimes_{B} A^{\prime \prime} \xrightarrow{\sim} M^{\prime \prime}
$$

as desired. (A similar argument for the case of general $M^{\prime}$ is given in [4, §1, Proposition 2].)

Corollary 3.6. With the notations as above, let $L$ be a $B$ module which may be inserted in a commutative diagram\\
\includegraphics[max width=\textwidth, alt={}, center]{133f4430-1007-4a22-9674-9093b9475bf3-10_401_424_1149_622}\\
where $q^{\prime}$ induces $L \otimes_{B} A^{\prime} \xrightarrow{\sim} M^{\prime}$. Then the canonical morphism $q^{\prime} \times q^{\prime \prime}: L \rightarrow N$ $=M^{\prime} \times_{M} M^{\prime \prime}$ is an isomorphism.

Proof. Apply Lemma 3.3 to the morphism $u=q^{\prime} \times q^{\prime \prime}$.\\
Remark. Lemma 3.4 is false, in general, if neither $A^{\prime \prime} \rightarrow A$ nor $A^{\prime} \rightarrow A$ is assumed surjective. For example, let $A^{\prime}$ be a sublocal ring of the local ring $A$, and map $A_{1}=A^{\prime \prime}$ into $A$ by inclusion. Let $a$ be a unit of $A$ such that the ideal $\left(a A^{\prime}\right) \cap A^{\prime}$ of $A^{\prime}$ is not flat ( = free) over $A^{\prime}$. (In $C_{\Lambda}$ one could take $A=k[t] /\left(t^{3}\right), A^{\prime}=k\left[t^{2}\right]$, $a=1+t$.) Let $M^{\prime}=M^{\prime \prime}=A^{\prime}, \quad M=A, \quad u^{\prime}=$ inclusion, $u^{\prime \prime}=$ multiplication by $a^{-1}$. Then $B \cong A^{\prime}$, while $N \cong\left(a A^{\prime}\right) \cap A^{\prime}$ is not flat over $B$.\\
Proof of Proposition 3.2. Let $u^{\prime}:\left(A^{\prime}, \eta^{\prime}\right) \rightarrow(A, \eta), u^{\prime \prime}:\left(A^{\prime \prime}, \eta^{\prime \prime}\right) \rightarrow(A, \eta)$ be morphisms of couples, where $u^{\prime \prime}$ is a surjection. Let $L^{\prime}, L, L^{\prime \prime}$ be corresponding invertible sheaves on $X^{\prime}=X_{A^{\prime}}, Y=X_{A}$, and $X^{\prime \prime}=X_{A^{\prime \prime}}$. Then we have morphisms $p^{\prime}: L^{\prime} \rightarrow L$, $p^{\prime \prime}: L^{\prime \prime} \rightarrow L$ (of sheaves on the topological space $\left|X_{0}\right|$, compatible with $\mathfrak{O}_{X^{\prime}} \rightarrow \mathfrak{D}_{Y}$, $\mathfrak{O}_{X^{\prime \prime}} \rightarrow \mathfrak{D}_{Y}$ ) which induce isomorphisms $L^{\prime} \otimes_{A^{\prime}} A \xrightarrow{\sim} L, L^{\prime \prime} \otimes_{A^{\prime \prime}} A \xrightarrow{\sim} L$.

Let $B=A^{\prime} \times{ }_{A} A^{\prime \prime}$, and let $Z=X_{B}$. Then we have a commutative diagram\\
\includegraphics[max width=\textwidth, alt={}, center]{133f4430-1007-4a22-9674-9093b9475bf3-11_429_439_401_624}\\
of sheaves on $\left|X_{0}\right|$; thus by Corollary 3.6 there is a canonical isomorphism $\mathfrak{O}_{Z} \xrightarrow{\sim} \mathfrak{D}_{X^{\prime}} \times_{\mathfrak{O}_{Y}} \mathfrak{O}_{X^{\prime \prime}}$, where $\mathfrak{O}_{X^{\prime}} \times_{\mathfrak{O}_{Y}} \mathfrak{O}_{X^{\prime \prime}}$ is the sheaf of $B$-algebras whose sections over an open $U$ in $\left|X_{0}\right|$ are given by

$$
\mathfrak{O}_{X^{\prime}} \times \mathfrak{D}_{Y} \mathfrak{O}_{X^{\prime \prime}}(U)=\mathfrak{O}_{X^{\prime}}(U) \times_{\mathfrak{D}_{Y}(U)} \mathfrak{O}_{X^{\prime \prime}}(U) .
$$

Hence $N=L^{\prime} \times_{L} L^{\prime \prime}$ is a sheaf on $Z$, obviously invertible, and the projections of $N$ on $L^{\prime}$ and $L^{\prime \prime}$ induce isomorphisms $N \otimes_{B} A^{\prime} \xrightarrow{\sim} L^{\prime}, N \otimes_{B} A^{\prime \prime} \xrightarrow{\sim} L^{\prime \prime}$ by Lemma 3.4.\\
If $M$ is another invertible sheaf on $Z$ for which there exist isomorphisms

$$
M \otimes A^{\prime} \xrightarrow{\sim} L^{\prime}, \quad M \otimes A^{\prime \prime} \xrightarrow{\sim} L^{\prime \prime},
$$

we have morphisms $q^{\prime}: M \rightarrow L^{\prime}, q^{\prime \prime}: M \rightarrow L^{\prime \prime}$ which induce these isomorphisms, and thus a commutative diagram\\
\includegraphics[max width=\textwidth, alt={}, center]{133f4430-1007-4a22-9674-9093b9475bf3-11_248_436_1525_631}\\
Here $\theta$ is the automorphism of $L$ given by the composition

$$
L \xrightarrow{\sim} L^{\prime} \otimes_{A^{\prime}} A \xrightarrow{\sim} M \otimes_{B} A \xrightarrow{\sim} L^{\prime \prime} \otimes_{A^{\prime}} A \xrightarrow{\sim} L .
$$

By hypothesis (ii) of 3.2, $\theta$ is multiplication by some unit $a \in A$. Lifting $a$ back to $a^{\prime \prime}$ in $A^{\prime \prime}$, we can change $q^{\prime \prime}$ to $a^{\prime \prime} q^{\prime \prime}$; thus we may assume that $u^{\prime} q^{\prime}=u^{\prime \prime} q^{\prime \prime}$. It follows from Corollary 3.6 that $M \xrightarrow{\sim} N$. We have therefore proved that

$$
P\left(A^{\prime} \times{ }_{A} A^{\prime \prime}\right) \xrightarrow{\sim} P\left(A^{\prime}\right) \times_{P(A)} P\left(A^{\prime \prime}\right)
$$

for any surjection $A^{\prime \prime} \rightarrow A$ in $C$.

Finally, letting $Y=X_{k[\varepsilon]}$, we have $\mathfrak{O}_{Y}=\mathfrak{D}_{X_{0}} \oplus \varepsilon \mathfrak{D}_{X_{0}}$, so there is a split exact sequence

$$
0 \longrightarrow \mathfrak{D}_{X_{0}} \xrightarrow{\exp } \mathfrak{D}_{Y}^{*} \longrightarrow \mathfrak{D}_{X_{0}}^{*} \longrightarrow 1
$$

where exp maps the (additive) sheaf $\mathfrak{O}_{X_{0}}$ into $\mathfrak{O}_{Y}^{*}$ by $\exp (f)=1+\varepsilon f$. Hence

$$
F(k[\varepsilon]) \cong \operatorname{ker}\left\{H^{1}\left(X_{0}, \mathfrak{O}_{Y}^{*}\right) \rightarrow H^{1}\left(X_{0}, \mathfrak{O}_{X_{0}}^{*}\right)\right\} \cong H^{1}\left(X_{0}, \mathfrak{O}_{X_{0}}\right)
$$

which has finite dimension, by assumption. This completes the proof of Proposition 3.2.\\
(3.7) Formal moduli. Let $X$ be a fixed prescheme over $k$, and $A \in C$. By an (infinitesimal) deformation of $X / k$ to $A$ we mean a product diagram\\
\includegraphics[max width=\textwidth, alt={}, center]{133f4430-1007-4a22-9674-9093b9475bf3-12_209_650_840_513}\\
where $Y$ is flat over $\operatorname{Spec} A$ and $i$ is (necessarily) a closed immersion. We will suppress the $i$ and refer to $Y$ as a deformation, if no confusion is possible. If $Y^{\prime}$ is another deformation to $A$ then $Y$ and $Y^{\prime}$ are isomorphic if there exists a morphism $f: Y \rightarrow Y^{\prime}$ over $A$ which induces the identity on the closed fibre $X$. ( $f$ must then be an isomorphism of preschemes, by Lemma 3.3.) Given the deformation $Y$ over $A$ and a morphism $A \rightarrow B$ in $C$, one has evidently an induced deformation $Y \otimes_{A} B$ over $B$; and if $Z$ is a deformation over $B$, one can define the notion of morphism $Z \rightarrow Y$ of deformations. (Notice that there is a one-to-one correspondence between such morphisms and the isomorphisms $Z \xrightarrow{\sim} Y \otimes_{A} B$ which they induce.

Define the deformation functor $\boldsymbol{D}=\boldsymbol{D}_{\boldsymbol{X} / \boldsymbol{k}}$ by setting

$$
D(A)=\text { Set of isomorphism classes of deformations of } X / k \text { to } A \text {. }
$$

We shall find that, in general, $D$ is not pro-representable, but that with rather weak finiteness restrictions on $X, D$ will have a hull.

Suppose that $\left(A^{\prime}, \eta^{\prime}\right) \rightarrow(A, \eta)$ and $\left(A^{\prime \prime}, \eta^{\prime \prime}\right) \rightarrow(A, \eta)$ are morphisms of couples, where $A^{\prime \prime} \rightarrow A$ is a surjection. Letting $X^{\prime}, Y, X^{\prime \prime}$ denote deformations in the class of $\eta^{\prime}, \eta, \eta^{\prime \prime}$ respectively, we have a diagram\\
\includegraphics[max width=\textwidth, alt={}, center]{133f4430-1007-4a22-9674-9093b9475bf3-12_173_316_1859_672}\\
of deformations. Therefore we can construct, as in the proof of 3.2 the sheaf $\mathfrak{O}_{X^{\prime}} \times_{\mathfrak{D}_{Y}} \mathfrak{O}_{X^{\prime \prime}}$ of $A^{\prime} \times_{A} A^{\prime \prime}$ algebras, and $\left(|X|, \mathfrak{O}_{X^{\prime}} \times_{\mathfrak{D}_{Y}} \mathfrak{O}_{X^{\prime \prime}}\right)$ defines a prescheme $Z$ flat over $A^{\prime} \times{ }_{A} A^{\prime \prime}$. (The fact that $Z$ is actually a prescheme consists of straight-forward checking; in fact it is the sum of $X^{\prime}$ and $X^{\prime \prime}$ in the category of preschemes\\
under $Y$, homeomorphic to $Y . Z$ is flat over $A^{\prime} \times_{A} A^{\prime \prime}$ by Lemma 3.4.) Furthérmore the closed immersions $X \rightarrow Y \rightarrow Z$ give $Z$ a structure of deformation of $X / k$ to $A^{\prime} \times{ }_{A} A^{\prime \prime}$ such that\\
\includegraphics[max width=\textwidth, alt={}, center]{133f4430-1007-4a22-9674-9093b9475bf3-13_396_408_470_626}\\
is a commutative diagram of deformations. In particular this shows that

$$
\boldsymbol{D}\left(A^{\prime} \times{ }_{A} A^{\prime}\right) \rightarrow \boldsymbol{D}\left(A^{\prime}\right) \times_{\boldsymbol{D}(A)} \boldsymbol{D}\left(A^{\prime}\right)
$$

is surjective, for every surjection $A^{\prime \prime} \rightarrow A$. That is, condition ( $\mathrm{H}_{1}$ ) of 2.11 is satisfied.\\
Suppose now that $W$ is another deformation over $B$, inducing the deformations\\
\includegraphics[max width=\textwidth, alt={}, center]{133f4430-1007-4a22-9674-9093b9475bf3-13_244_432_1145_611}\\
$X^{\prime}$ and $X^{\prime \prime}$. Then there is a commutative diagram of deformations, where $\theta$ is the composition

$$
Y \xrightarrow{\sim} X^{\prime} \otimes_{A^{\prime}} A \xrightarrow{\sim} W \otimes_{B} A \longrightarrow X^{\prime \prime} \otimes_{A^{\prime \prime}} A \xrightarrow{\sim} Y .
$$

If $\theta$ can be lifted to an automorphism $\theta^{\prime}$ of $X^{\prime}$, such that $\theta^{\prime} u^{\prime}=u^{\prime} \theta$, then we can replace $q^{\prime}$ with $q^{\prime} \theta^{\prime}$; then we would have an isomorphism $W \xrightarrow{\sim} Z$ by Corollary 3.6. Now if $A=k$ (so that $Y=X, \theta=\mathrm{id}$ ) $\theta^{\prime}$ certainly exists, so condition ( $\mathrm{H}_{2}$ ) is satisfied.

To consider the condition $\left(\mathrm{H}_{4}\right)$, let $p:\left(A^{\prime}, \eta^{\prime}\right) \rightarrow(A, \eta)$ be a morphism of couples, where $p$ is a small extension. For each morphism $B \rightarrow A$, let $D_{\eta}(B)$ denote as usual the set of $\zeta \in D(B)$ such that $\zeta \otimes_{B} A=\eta$. Pick a deformation $Y^{\prime}$ in the class of $\eta^{\prime}$; then

Lemma 3.8. The following are equivalent\\
(i) $\boldsymbol{D}_{\eta}\left(A^{\prime} \times{ }_{A} A^{\prime}\right) \xrightarrow{\sim} \boldsymbol{D}_{\eta}\left(A^{\prime}\right) \times \boldsymbol{D}_{\eta}\left(A^{\prime}\right)$,\\
(ii) Every automorphism of the deformation $Y=Y^{\prime} \otimes_{A^{\prime}} A$ is induced by an automorphism of the deformation $Y^{\prime}$.

Proof. (i) ⇒ (ii). Let $u: Y \rightarrow Y^{\prime}$ be the induced morphism of deformations.

If $\theta$ is an automorphism of $Y$, then one can construct deformations $Z, W$ over $A^{\prime} \times_{A} A^{\prime}$ to yield "sum diagrams"\\
\includegraphics[max width=\textwidth, alt={}, center]{133f4430-1007-4a22-9674-9093b9475bf3-14_399_400_430_360}\\
\includegraphics[max width=\textwidth, alt={}, center]{133f4430-1007-4a22-9674-9093b9475bf3-14_400_385_429_901}\\
of deformations. Since $Z$ and $W$ have isomorphic projections on both factors, there is an isomorphism $\rho: Z \xrightarrow{\sim} W . \rho$ induces automorphisms $\theta_{1}$ and $\theta_{2}$ of $Y^{\prime}$, and an automorphism $\phi$ of $Y$ such that

$$
\theta_{1} u \theta=u \phi, \quad \theta_{2} u=u \phi .
$$

Therefore $u \theta=\theta_{1}^{-1} \theta_{2} u$ and $\theta_{1}^{-1} \theta_{2}$ induces $\theta$.\\
(ii) ⇒ (i). In a similar manner, it follows from (ii) that $t_{F} \otimes I(I=\operatorname{ker} p)$ acts freely on $\eta^{\prime}$ (i.e., $\left(\eta^{\prime}\right)^{\sigma}=\eta^{\prime}$ implies $\sigma=0$ ). Since the action of $t_{F} \otimes I$ on $\boldsymbol{D}_{\eta}\left(A^{\prime}\right)$ is transitive, it follows that $D_{\eta}\left(A^{\prime}\right)$ is a principal homogeneous space under $t_{F} \otimes I$, which is equivalent to (i).

It should be remarked that the obstruction to lifting $\theta$ lies in $t_{F} \otimes I$ and is often nonzero (see e.g., [4, §4]).

Finally it remains to consider the finiteness condition $\left(\mathrm{H}_{3}\right)$. If $X$ is smooth over $k$ (in ancient terminology absolutely simple), then Grothendieck has shown in S.G.A. III, Theorem 6.3, that

$$
t_{D} \cong H^{1}(X, \Theta)
$$

where $\Theta$ is the tangent sheaf of $X$ over $k$. Thus $t_{D}$ has finite dimension if $X$ is smooth and proper over $k$. In general, it is shown in [4] that for any scheme $X$ locally of finite type over $k$, there is an exact sequence


\begin{equation*}
0 \rightarrow H^{1}\left(X, T^{0}\right) \rightarrow t_{D} \rightarrow H^{0}\left(X, T^{1}\right) \rightarrow H^{2}\left(X, T^{0}\right) \tag{3.9}
\end{equation*}


where $T^{0}$ is the sheaf of derivations of $\mathfrak{O}_{X}$, and $T^{1}$ is a (coherent) sheaf isomorphic to the sheaf of germs of deformations of $X / k$ to $k[\varepsilon]$. If $X$ is smooth over $k$, then $T^{0}=\Theta, T^{1}=0$. Thus, in summary

Proposition 3.10. If $X$ is either\\
(a) proper over $k$ or\\
(b) affine with only isolated singularities, then $\boldsymbol{D}$ has a hull $(R, \xi)$. $(R, \xi)$ pro-represents $\boldsymbol{D}$ if and only if for each small extension $A^{\prime} \rightarrow A$, and each deformation $Y^{\prime}$ of $X / k$ to $A^{\prime}$, every automorphism of the deformation $Y^{\prime} \otimes_{A^{\prime}} A$ is induced by an automorphism of $Y^{\prime}$.\\
(3.11) The automorphism functor. One can formalize the obstructions to prorepresenting $D$ as follows. Let $X$ be a prescheme proper over $k$, and let $(R, \xi)$ be a hull of the deformation functor $D . \xi$ is represented by a formal prescheme $\mathfrak{X}=\operatorname{inj} \operatorname{Lim} X_{n}$ over $R$, where $X_{n}$ is a deformation of $X / k$ to $R / \boldsymbol{m}^{n}$. For each morphism $R \rightarrow A$ in $C_{\Lambda}$, we get a deformation $\mathfrak{X}_{A}=\mathfrak{X} \times_{\operatorname{Spec} R} \operatorname{Spec} A$ of $X / k$ to $A$. We can therefore define a group functor $A$ on the category $C_{R}$ of Artin local $R$-algebras:

$$
A: A \mapsto \text { group of automorphisms of the deformation } \mathfrak{X}_{A} .
$$

If $A^{\prime} \rightarrow A$ and $A^{\prime \prime} \rightarrow A$ are morphisms in $C_{R}$ with $A^{\prime \prime} \rightarrow A$ a surjection, and if we put $B=A^{\prime} \times{ }_{A} A^{\prime \prime}$ then we have a canonical isomorphism, respecting the structures as deformations:

$$
\mathfrak{O}_{x_{B}} \cong \mathfrak{D}_{X_{A^{\prime}}} \times \mathfrak{D}_{X_{A}} \mathfrak{O}_{X_{A}} .
$$

by Corollary 3.6. It follows easily that (2.12) is an isomorphism, so that $\left(\mathrm{H}_{1}\right)$, $\left(\mathrm{H}_{2}\right)$ and $\left(\mathrm{H}_{4}\right)$ of Theorem 2.11 are satisfied. Finally the computations of Grothendieck in S.G.A. III, §6, show that the tangent space of $A$ is given by

$$
t_{A / R} \cong H^{0}\left(X_{0}, T^{0}\right)
$$

where $T^{0}$ is, again, the (coherent) sheaf of derivations of $\mathfrak{O}_{X}$ over $k$. Thus $t_{A}$ has finite dimension, and we find:

Proposition 3.12. If $X$ is proper over $k$, the functor $A$ is pro-represented by a complete local $R$ algebra, $S$, which is a group object in the category dual to $\hat{C}_{R}$ (i.e., $S$ is a formal Lie group over $R$ ). The deformation functor $D$ is pro-representable (by $R$ ) if and only if $S$ is a power series ring over $R$.

The last statement follows from Lemma 3.8 and the smoothness criterion of Remark 2.10.

In a future paper I will discuss the deformation functor in more detail, with particular attention to the contribution of singular points on $X$.

\section*{References}
\begin{enumerate}
  \item N. Bourbaki, Algèbre commutative, Chapitre III, Actualités Sci. Ind., 1923.
  \item A. Grothendieck, Séminaire de géomètrie algébrique (S.G.A.), Inst. Hautes Etudes Sci., Paris.
  \item -, Technique de descent et théoremès d'existence en géomètrie algébrique, II, Séminaire Bourbaki, Exposé 195, 1959/1960.
  \item A. H. M. Levelt, Sur la proreprésentabilité de certains foncteurs en géomètrie algébrique, Notes, Katholieke Universiteit, Nijmegen, Netherlands.
  \item M. Schlessinger, Infinitesimal deformations of singularities, Ph.D. Thesis, Harvard Univ., Cambridge, Mass., 1964.
\end{enumerate}


\end{document}